\documentclass[12pt, a4paper]{article}

%=============== PAQUETES ===============
\usepackage[utf8]{inputenc}
\usepackage[T1]{fontenc}
\usepackage[spanish,es-tabla]{babel}
\usepackage{circuitikz}
\usepackage{siunitx}
\usepackage{geometry}

%=============== CONFIGURACIÓN ===============
\geometry{a4paper, margin=2.5cm}

\sisetup{
    output-decimal-marker = {.},
    per-mode = symbol
}

\begin{document}

\begin{figure}[h!]
    \centering
    \begin{circuitikz}[scale=1]
        \draw (0,3) node[anchor=east] {A} to[short, o-*] (1,3)
              to[R, l=$R_1$ (\SI{5}{\ohm})] (4,3)
              to[short, *-] (4,4)
              to[R, l=$R_2$ (\SI{20}{\ohm})] (7,4)
              to[short, -*] (7,3);
        \draw (4,3) to[short] (4,2)
              to[R, l=$R_3$ (\SI{20}{\ohm})] (7,2)
              to[short] (7,3);
        \draw (1,3) to[short] (1,1)
              to[R, l=$R_4$ (\SI{15}{\ohm})] (7,1)
              to[short] (7,3);
        \draw (7,3) to[short] (8,3)
              to[R, l=$R_5$ (\SI{2.5}{\ohm})] (8,0)
              to[short, -o] (0,0) node[anchor=east] {B};
    \end{circuitikz}
    \caption{Circuito para el problema de reducción de resistencias.}
    \label{fig:ej2}
\end{figure}

\end{document}